Se logró determinar la constante de Planck a partir de las tres líneas de la serie de Balmer,
obteniendo un valor de \( h = \planckhexp \), con un error relativo porcentual de \qty{0.126}{\percent}
respecto al valor reportado en la literatura.

En cuanto al elemento desconocido analizado, al observar su serie espectral en el rango visible y  que la
segunda línea medida presenta la mayor intensidad y un color amarillo, se concluye que es sodio.
